\documentclass[10pt,a4paper]{amsart}
\usepackage[margin=19mm]{geometry}
\usepackage{amsmath,amssymb,mathtools}
\usepackage[scr=rsfs,cal=euler]{mathalfa}
\usepackage{xfrac}
\usepackage[unicode,hidelinks,bookmarks=false]{hyperref}
\newcommand{\ie}{i.e.\ }
\newcommand{\cf}{cf.\ }
\newcommand{\resp}{respectively\ }
\newcommand{\Subsection}{\subsection}
\providecommand{\nobreakdash}{\nobreak\hbox{-}}
\setlength{\emergencystretch}{2em}
\multlinegap0pt
\usepackage{bm}
\usepackage{bbm}
\usepackage{dsfont}
\usepackage{xspace}
\usepackage{cancel}
\usepackage{mathtools}
\usepackage{amsthm}
\makeatletter
\@ifundefined{proposition}{\newtheorem{proposition}{Proposition}}{}
\@ifundefined{theorem}{\newtheorem{theorem}{Theorem}}{}
\makeatother
\def \geq{\geqslant}
\def \leq{\leqslant}
\def\R {\mathbb{R}}
\title[Small-time approximate controllability for the nonlinear com]{Small-time approximate controllability for the nonlinear complex Ginzburg-Landau equation with bilinear control}
\author{Xingwu Zeng, Can Zhang}
\date{Source: arXiv:2512.22512v1 (CC BY 4.0)}
\begin{document}
\maketitle

\noindent\emph{Source and licence.} Small-time approximate controllability for the nonlinear complex Ginzburg-Landau equation with bilinear control. Version arXiv:2512.22512v1 (CC BY 4.0), distributed under the Creative Commons Attribution 4.0 International licence, as recorded in the arXiv licence metadata for this exact version and shown on the abstract page https://arxiv.org/abs/2512.22512v1. Full rights notice: licenses/arxiv-2512.22512-CC-BY-4.0.txt.\par\smallskip

\noindent\textit{Editorial scope.} Counted unit: the Theorem whose source label is \texttt{thm2.2} in the article (source0.tex lines 442--448, source label \texttt{thm2.2}), delivered together with the article's own complete original proof of it (source0.tex lines 449--522, one \texttt{proof} environment of 74 lines). No mathematics is written, repaired or re-derived: statement and proof are the article's own text line for line. Required context retained uncounted: (0.2) (lines 100--102); (0.1) (lines 199--201); prop1.1 (lines 242--257); stability (lines 249--251); prop1.2 (lines 269--275). Every definition, display and label of those lines is retained unchanged. Omitted from this package and not redistributed: 1--99, 103--198, 202--241, 252--268, 276--441, 523--918. Nothing inside a retained excerpt is omitted or altered: every retained body is delivered exactly as the article's lines are written, so no omission receipt of any kind accompanies this package. Citations: the retained excerpts carry no citation command at all, so no bibliography entry of the article is transcribed and no reference is redistributed. Reference adaptation: none was needed and none was made; every reference inside the delivered unit resolves inside it, so no unresolved reference or citation is produced. Printed slips kept literal: no printed word, symbol, hypothesis or number of the counted statement or of its proof was altered. Layout and class: the article's own class is \texttt{elsarticle} with packages including lineno, hyperref, mathrsfs, amsmath, amssymb, amsthm, color, esint, enumitem, CJK; the wrapper is the lane's pinned \texttt{amsart} editorial wrapper at 10pt on A4, which restates the theorem-like environments the retained text needs and adds the article's own macro definitions that the retained text needs (\texttt{\textbackslash R}, \texttt{\textbackslash geq}, \texttt{\textbackslash leq}), with extra wrapper packages (bm, bbm, dsfont, xspace, cancel, mathtools). The delivered numbering is the unit's own. Assets: no retained span contains a figure environment, a graphics-inclusion command, a PDF-page inclusion, a table or a TikZ picture; no image file is copied into this package, the package declares no asset dependency and no image file is redistributed with it. Licence: version arXiv:2512.22512v1 of this article is distributed under the Creative Commons Attribution 4.0 International licence, as recorded in the arXiv licence metadata for this exact version and shown on the abstract page https://arxiv.org/abs/2512.22512v1. A full rights notice is delivered at licenses/arxiv-2512.22512-CC-BY-4.0.txt. Characters: every retained excerpt is pure ASCII, so the delivered engine, XeLaTeX with its default Latin Modern text font, reports no missing character.
\begin{equation}\label{(0.2)}
\psi(0, x)=\psi_0(x)
\end{equation}

\begin{equation}\label{(0.1)}
	\partial_t \psi=V \psi+(1+i \nu) \Delta \psi-(1+i \mu)|\psi|^{2 \sigma} \psi +(r_1+i r_2) \langle u(t), Q(x)\rangle \psi ,
\end{equation}

\begin{proposition}\label{prop1.1}
For any $s>d / 2,\, \tilde{\psi}_0 \in H^s$, and $\tilde{u} \in L_{\text{loc}}^2\left(\mathbb{R}_{+} ; \mathbb{R}^q\right)$, there is a maximal time $\tilde{T}=\tilde{T}\left(\tilde{\psi}_0, \tilde{u}\right)>0$ and a unique solution $\tilde{\psi}$ of the problem \eqref{(0.1)}, \eqref{(0.2)} with $\left(\psi_0, u\right)=\left(\tilde{\psi}_0, \tilde{u}\right)$, whose restriction to the interval $J_T$ belongs to $C\left(J_T ; H^s\right)$ for any $T<\tilde{T}$. If $\tilde{T}<\infty$, then $\|\tilde{\psi}(t)\|_s \rightarrow \infty$ as $t \rightarrow \tilde{T}^{-}$. Furthermore, for any $T<\tilde{T}$, there are constants $\delta=\delta(T, \Lambda)>0$ and $C=C(T, \Lambda)>0$, where
$$
\Lambda=\|\tilde{\psi}\|_{C\left(J_T ; H^s\right)}+\|\tilde{u}\|_{L^2\left(J_T ; \mathbb{R}^q\right)},
$$
such that the following properties hold:\\
(i) For any $\psi_0 \in H^s$ and $u \in L^2\left(J_T ; \mathbb{R}^q\right)$ satisfying
\begin{equation}\label{stability}
 \|\psi_0-\tilde{\psi}_0 \|_s+\|u-\tilde{u}\|_{L^2\left(J_T ; \mathbb{R}^q\right)}<\delta,
\end{equation}
the problem \eqref{(0.1)}, \eqref{(0.2)} has a unique solution $\psi \in C\left(J_T ; H^s\right)$.\\
(ii) Let $\mathcal{R}$ be the resolving operator for \eqref{(0.1)}, i.e., the mapping taking a couple $\left(\psi_0, u\right)$ satisfying \eqref{stability} to the solution $\psi$. Then there exists a constant $C>0$ such that
\begin{equation}\label{continuity}
\left\|\mathcal{R}\left(\psi_0, u\right)-\mathcal{R}\left(\tilde{\psi}_0, \tilde{u}\right)\right\|_{C\left(J_T ; H^s\right)} \leq C\left( \|\psi_0-\tilde{\psi}_0 \|_s+\|u-\tilde{u}\|_{L^2\left(J_T ; \mathbb{R}^q\right)}\right) .
\end{equation}
\end{proposition}

\begin{proposition}\label{prop1.2}
For any integer $s \geq s_d,\, \psi_0 \in H^s,\, (u,r_1,r_2) \in \mathbb{R}^q\times\R\times\R$, and  non-negative function $\varphi \in C^r\left(\mathbb{T}^d ; \mathbb{R}\right)$ with $r=s+2$, there is a constant $\delta_0>0$ such that, for any $\delta \in\left(0, \delta_0\right)$, we have\footnote{For any vector $u \in \mathbb{R}^q$, with a slight abuse of notation, we denote by the same letter the constant value function equals to $u$.} $\delta^{-1} u \in \Theta\left(e^{-(a+ib) \delta^{-1 / 2} \varphi} \psi_0, \delta\right)$ and
\begin{equation*}\label{(1.3)}
e^{ (a+ib)\delta^{-1 / 2} \varphi} \mathcal{R}_\delta\left(e^{-(a+ib) \delta^{-1 / 2} \varphi} \psi_0, (\delta^{-1} u,r_1,r_2)\right) \rightarrow e^{(r_1+i r_2)\left( \mathbb{B}(\varphi)+\langle u, Q\rangle\right)} \psi_0 \quad \text { in } H^s \text { as } \delta \rightarrow 0^{+}.
\end{equation*}
Here $\mathcal{R}_\delta$ is the restriction of the solution at time $t=\delta$.
\end{proposition}

\begin{theorem}\label{thm2.2}
Assume that the condition~{\hyperlink{H1}{\textbf{($\mathscr{P}$)}}} is satisfied, $ \mathbf{1}\in \operatorname{span}\left\{Q_j\mid j=1, \ldots, q\right\} $ and $ r_1=1,\, r_2=-\nu $.
Then for any $\varepsilon>0,\, \mathcal{T}>0,\, \psi_0 \in H^s$, where $s \geq s_d$ is an integer, and $\theta \in C^r\left(\mathbb{T}^d ; \mathbb{R}\right)$, there are $T \in(0, \mathcal{T})$ and $u \in \Theta\left(\psi_0, T\right) \cap C^{\infty}\left(J_T ; \mathbb{R}^q\right)$ such that
\begin{equation}\label{(2.2)}
\left\|\mathcal{R}_T\left(\psi_0, (u,1,-\nu)\right)-e^{(1-i\nu) \theta} \psi_0\right\|_s<\varepsilon.
\end{equation}
\end{theorem}

\begin{proof}
We first let $ r_1=a=\frac{1}{1+\nu^2},\, r_2 = b =\frac{-\nu}{1+\nu^2} $, and
use an induction argument in $N=0,1,2,...$ to show that the approximate controllability in this theorem is true for any $\theta \in H_N$. 
More precisely, we prove the following property:
\begin{description}\item[\hypertarget{PN}]{\textbf{(P$_N$)}}  
For any $\theta \in H_N$ and $\psi_0 \in H^s$, there is a family $\left\{u_\tau\right\}_{\tau>0} \subset L^2\left(J_1 ; \mathbb{R}^q\right)$ such that $u_\tau \in \Theta\left(\psi_0, \tau\right)$ for sufficiently small $\tau>0$, and
\begin{equation}\label{(2.3)}
\mathcal{R}_\tau\left(\psi_0, \left(u_\tau,\frac{1}{1+\nu^2},\frac{-\nu}{1+\nu^2}\right)\right) \rightarrow e^{\frac{1-i\nu}{1+\nu^2} \theta} \psi_0 \quad \text { in } H^s \text { as } \tau \rightarrow 0^{+}.
\end{equation}
\end{description} 
Combined with the saturation condition~{\hyperlink{H1}{\textbf{($\mathscr{P}$)}}}, this leads to the approximate controllability for any $\theta \in C^r\left(\mathbb{T}^d ; \mathbb{R}\right)$.
\\
\textbf{Step~1.} Case $N=0$. Applying Proposition~\ref{prop1.2} for $\varphi=0$ and $u \in \mathbb{R}^q$ with $\theta=$ $\langle u, Q\rangle$, we obtain
$$
\mathcal{R}_\delta\left(\psi_0, \left(\delta^{-1} u,\frac{1}{1+\nu^2},\frac{-\nu}{1+\nu^2}\right)\right) \rightarrow e^{\frac{1-i\nu}{1+\nu^2} \theta } \psi_0 \quad \text { in } H^s \text { as } \delta \rightarrow 0^{+}.
$$
This implies \eqref{(2.3)} with $\tau=\delta$ and $u_\tau=\delta^{-1} u$.
\\
\textbf{Step~2.} Case $N \geq 1$. We assume that \hyperlink{PN}{\textbf{(P$_{N-1}$)}} is true. Let $\tilde{\theta} \in H_N$ be of the form
$$
\tilde{\theta}=\theta_0+\sum_{j=1}^n \mathbb{B}\left(\theta_j\right),
$$
where $n \geq 1$ and $\theta_j \in H_{N-1},\, j=0, \ldots, n$.
Take $ \theta_1 $ and $ c>0 $ to be such that $ \tilde{\theta}_1=\theta_1 +c\geq0 $. Note that $ \mathbb{B}(\tilde{\theta}_1) = \mathbb{B}(\theta_1) $. Applying Proposition~\ref{prop1.2} with $\varphi=\tilde{\theta}_1$ and $u=0$, we get
\begin{equation*}\label{(2.4)}
e^{\frac{1-i\nu}{1+\nu^2}\delta^{-1 / 2} \tilde{\theta}_1} \mathcal{R}_\delta\left(e^{-\frac{1-i\nu}{1+\nu^2} \delta^{-1 / 2} \tilde{\theta}_1} \psi_0, \left(0,\frac{1}{1+\nu^2},\frac{-\nu}{1+\nu^2}\right)\right) \rightarrow e^{\frac{1-i\nu}{1+\nu^2} \mathbb{B}\left(\theta_1\right)} \psi_0 \quad \text { in } H^s \text { as } \delta \rightarrow 0^{+} .
\end{equation*}
Since $ c\in H_0$ and $\theta_1\in H_{N-1} $, we have $ \tilde{\theta}_1\in H_{N-1} $. 
The induction hypothesis implies that, for any $\delta>0$, there are families of controls $\left\{u_{\tau, \delta}^1\right\} \subset L^2\left(J_1 ; \mathbb{R}^q\right)$ and $\left\{u_{\tau, \delta}^2\right\} \subset L^2\left(J_1 ; \mathbb{R}^q\right)$ such that $u_{\tau, \delta}^1 \in \Theta\left(\psi_0, \tau\right)$ and $u_{\tau, \delta}^2 \in \Theta\left(\mathcal{R}_\delta\left(e^{-\frac{1-i\nu}{1+\nu^2} \delta^{-1 / 2} \tilde{\theta}_1} \psi_0, \left(0,\frac{1}{1+\nu^2},\frac{-\nu}{1+\nu^2}\right)\right), \tau\right)$ for sufficiently small $\tau>0$, and
\begin{equation*}\label{(2.5)}
\mathcal{R}_\tau\left(\psi_0, \left(u_{\tau, \delta}^1,\frac{1}{1+\nu^2},\frac{-\nu}{1+\nu^2}\right)\right) \rightarrow e^{-\frac{1-i\nu}{1+\nu^2} \delta^{-1 / 2} \tilde{\theta}_1} \psi_0,
\end{equation*}
\begin{align}
&\mathcal{R}_\tau\left(\mathcal{R}_\delta\left(e^{-\frac{1-i\nu}{1+\nu^2} \delta^{-1 / 2} \tilde{\theta}_1} \psi_0, \left(0,\frac{1}{1+\nu^2},\frac{-\nu}{1+\nu^2}\right)\right), \left(u_{\tau, \delta}^2,\frac{1}{1+\nu^2},\frac{-\nu}{1+\nu^2}\right)\right) \nonumber\\
 &\rightarrow e^{\frac{1-i\nu}{1+\nu^2} \delta^{-1 / 2} \tilde{\theta}_1} \mathcal{R}_\delta\left(e^{-\frac{1-i\nu}{1+\nu^2} \delta^{-1 / 2} \tilde{\theta}_1} \psi_0, \left(0,\frac{1}{1+\nu^2},\frac{-\nu}{1+\nu^2}\right)\right) \nonumber\label{(2.6)}
\end{align}
in $H^s$ as $\tau \rightarrow 0^{+}$. 

Combining these controls and using Proposition~\ref{prop1.1}, we can construct a new family $\left\{u_\tau^1\right\} \subset L^2\left(J_1 ; \mathbb{R}^q\right)$ such that $u_\tau^1 \in \Theta\left(\psi_0, \tau\right)$ for sufficiently small $\tau>0$, and
$$
\mathcal{R}_\tau\left(\psi_0, \left(u_\tau^1,\frac{1}{1+\nu^2},\frac{-\nu}{1+\nu^2}\right)\right) \rightarrow e^{\frac{1-i\nu}{1+\nu^2} \mathbb{B} (\theta_1 )} \psi_0 \quad \text { in } H^s \text { as } \tau \rightarrow 0^{+} .
$$
Iterating this argument with $\theta_j \in H_{N-1},\, j=0, \ldots, n$, we obtain a family $\left\{u_\tau^n\right\} \subset$ $L^2\left(J_1 ; \mathbb{R}^q\right)$ such that $u_\tau^n \in \Theta\left(\psi_0, \tau\right)$ for small $\tau>0$ and
$$
\mathcal{R}_\tau\left(\psi_0, \left(u_\tau^n,\frac{1}{1+\nu^2},\frac{-\nu}{1+\nu^2}\right)\right) \rightarrow e^{\frac{1-i\nu}{1+\nu^2}\left(\theta_0+\sum_{j=1}^n \mathbb{B} (\theta_j )\right)} \psi_0=e^{\frac{1-i\nu}{1+\nu^2} \tilde{\theta}} \psi_0 \quad \text { in } H^s \text { as } \tau \rightarrow 0^{+} .
$$
As $\tilde{\theta} \in H_N$ is arbitrary, this showes the required Property~\hyperlink{PN}{\textbf{(P$_{N}$)}} for $N$.
\\
\textbf{Step~3. Conclusion.}  Finally, let $\theta \in C^r\left(\mathbb{T}^d ; \mathbb{R}\right)$ be arbitrary. It is clear that $ (1+\nu^2)\theta\in C^r\left(\mathbb{T}^d ; \mathbb{R}\right) $. By the saturation hypothesis~{\hyperlink{H1}{\textbf{($\mathscr{P}$)}}}, $H_{\infty}$ is dense in $C^r\left(\mathbb{T}^d ; \mathbb{R}\right)$. This implies that we can find $N \geq 1$ and $\tilde{\theta} \in H_N$ such that
$$
\left\|e^{(1-i\nu)\theta} \psi_0-e^{\frac{1-i\nu}{1+\nu^2} \tilde{\theta}} \psi_0\right\|_s<\varepsilon.
$$
Applying \hyperlink{PN}{\textbf{(P$_{N}$)}}  for $\tilde{\theta} \in H_N$, we find $T \in(0, \mathcal{T})$ and $\tilde{u} \in \Theta\left(\psi_0, T\right)$ such that \begin{equation*}
	\left\|\mathcal{R}_T\left(\psi_0, \left(\tilde{u},\frac{1}{1+\nu^2},\frac{-\nu}{1+\nu^2}\right)\right)-e^{(1-i\nu) \theta} \psi_0\right\|_s<2\varepsilon.
\end{equation*}
Since
\begin{equation*}
\mathcal{R}_T\left(\psi_0, \left(\tilde{u},\frac{1}{1+\nu^2},\frac{-\nu}{1+\nu^2}\right)\right) = \mathcal{R}_T\left(\psi_0, \left(\frac{\tilde{u}}{1+\nu^2},1,-\nu\right)\right),
\end{equation*}
we obtain \eqref{(2.2)} with $ u= \frac{\tilde{u}}{1+\nu^2} $.

Proposition~\ref{prop1.1} and a density argument show that we can take $u \in \Theta\left(\psi_0, T\right) \cap C^{\infty}\left(J_T ; \mathbb{R}^q\right)$. Indeed, since $C^{\infty}\left(J_T ; \mathbb{R}^q\right)$ is dense in $L^2\left(J_T ; \mathbb{R}^q\right)$, for any $u\in  \Theta\left(\psi_0, T\right) $ such that  \eqref{(2.2)} is satisfied, one can choose $u_\varepsilon\in C^{\infty}\left(J_T ; \mathbb{R}^q\right)$ satisfying
\begin{equation*}
	\|u_\varepsilon - u\|_{L^2\left(J_T ; \mathbb{R}^q\right)}\leq \varepsilon,
\end{equation*}
this, along with \eqref{(2.2)} and  Proposition~\ref{prop1.1}, finally implies that $u_\varepsilon\in  \Theta\left(\psi_0, T\right) $ and
\begin{equation*}
	\begin{aligned}
	&\left\|\mathcal{R}_T\left(\psi_0, (u_\varepsilon,1,-\nu)\right)-e^{(1-i\nu) \theta} \psi_0\right\|_s \\
	\leq 	&\left\|\mathcal{R}_T\left(\psi_0, (u_\varepsilon,1,-\nu)\right)- \mathcal{R}_T\left(\psi_0, (u,1,-\nu)\right)\right\|_s + \left\|\mathcal{R}_T\left(\psi_0, (u,1,-\nu)\right)-e^{(1-i\nu) \theta} \psi_0\right\|_s\\
	< &( C+1 )\varepsilon.
	\end{aligned}
\end{equation*}
\end{proof}

\end{document}
